% !TeX root = ../../Book2.tex
\section{线性系统、矩阵指数与变常数}
\label{b2:sec:ode-linear-systems}

线性系统的时间系数即使只有连续性，也能由积分迭代构造唯一解。这里的估计含有阶乘，不必把时间区间缩短到使积分算子成为压缩；后文 Sturm--Liouville 问题的通量系统正需要这个版本。矩阵指数与一般变常数公式仍留待撰写。

\subsection{连续系数线性系统与积分迭代}
\label{b2:subsec:070205}

\begin{proposition}\label{b2:prop:ode-continuous-linear}
设 \(n\) 为正整数，\(\mathbb K=\mathbb R\) 或 \(\mathbb C\)，\(\alpha<\beta\)，\(A\in C([\alpha,\beta];\mathbb K^{n\times n})\)。对任意 \(t_0\in[\alpha,\beta]\)、\(y_0\in\mathbb K^n\)，问题
\[
 Y'(t)=A(t)Y(t),\quad Y(t_0)=y_0
\]
在整个 \([\alpha,\beta]\) 上有唯一的 \(C^1\) 解；端点处的导数取相应单侧值。若 \(M=\max_{[\alpha,\beta]}\|A(t)\|\)，其中矩阵范数由 Euclid 范数诱导，则
\[
 |Y(t)|\leq e^{M|t-t_0|}\cdot|y_0|.
\]
\end{proposition}
\begin{proof}
令 \(Y_0(t)=y_0\)，递推定义
\[
 Y_{k+1}(t)=y_0+\int_{t_0}^t A(s)Y_k(s)\,\mathrm ds,
 \quad k\geq0.
\]
第一次差满足 \(|Y_1(t)-Y_0(t)|\leq M|t-t_0|\cdot|y_0|\)。反复积分并归纳得到
\[
 |Y_{k+1}(t)-Y_k(t)|
 \leq\frac{M^{k+1}|t-t_0|^{k+1}}{(k+1)!}\cdot|y_0|.
\]
当 \(t<t_0\) 时，把积分方向反转后使用同一估计。由于 \(|t-t_0|\leq\beta-\alpha\)，右端的数项级数收敛，故 \(Y_k\) 一致收敛到连续函数 \(Y\)。系数有界，允许在积分方程中取一致极限，得到
\[
 Y(t)=y_0+\int_{t_0}^t A(s)Y(s)\,\mathrm ds.
\]
被积函数连续，所以 \(Y\in C^1\) 并满足初值问题。把逐次差的估计求和，同时得到 \(|Y(t)|\leq e^{M|t-t_0|}\cdot|y_0|\)。

若 \(Y,\widetilde Y\) 是同初值的两解，令 \(Z=Y-\widetilde Y\)、\(D=\max_{[\alpha,\beta]}|Z|\)。由 \(Z(t)=\int_{t_0}^t A(s)Z(s)\,\mathrm ds\) 反复代入，对每个正整数 \(m\) 都有
\[
 |Z(t)|\leq\frac{M^m|t-t_0|^m}{m!}\cdot D.
\]
令 \(m\to\infty\) 得 \(Z(t)=0\)，故解唯一。
\end{proof}

矩阵初值问题可逐列应用本命题。若指定某时刻的初始矩阵，每一列唯一确定，因而整个矩阵解也唯一确定。证明只使用系数连续与紧区间上的有界性，没有对 \(A\) 求导。

\subsection{矩阵指数与变常数公式}
\label{b2:subsec:070206}

待撰写。拟先处理常系数矩阵指数，再建立基本解矩阵及非齐次系统的变常数公式。前一命题目前只给出连续系数齐次系统的初值理论。
